\textbf{结论名称：} 分层抽样中样本量的比例分配

\textbf{适用条件：} 总体由若干层组成，层间差异大、层内差异小；抽样时各层样本按比例分配，以保证样本结构与总体结构一致。

\textbf{变量说明：} 设总体容量为 \(N\)，分为 \(L\) 层，第 \(i\) 层的个体数为 \(N_i\)（\(i=1,2,\dots,L\)），满足 \(\sum_{i=1}^{L} N_i = N\)；样本容量为 \(n\)，从第 \(i\) 层抽取的样本量为 \(n_i\)。

\textbf{核心公式：} 
\[
\boxed{n_i = n \cdot \frac{N_i}{N}}
\]

\textbf{使用场景：} 大规模抽样调查（如人口普查、市场调查、教育测评），特别是各层个体数差异较大且关注各层代表性时，用于确定各层应抽取的样本数。

\textbf{简单例子：} 
某校高一共有学生 \(1200\) 人，其中男生 \(700\) 人，女生 \(500\) 人。现按性别分层，用比例分配抽取容量 \(n=120\) 的样本。则
男生抽取 \(n_1 = 120 \times \dfrac{700}{1200} = 70\) 人，女生抽取 \(n_2 = 120 \times \dfrac{500}{1200} = 50\) 人，合计 \(120\) 人，各层抽样比例均为 \(10\%\)。